\documentclass[10pt,a4paper]{amsart}
\usepackage[margin=19mm]{geometry}
\usepackage{amsmath,amssymb,mathtools}
\usepackage[scr=rsfs,cal=euler]{mathalfa}
\usepackage{xfrac}
\usepackage[unicode,hidelinks,bookmarks=false]{hyperref}
\newcommand{\ie}{i.e.\ }
\newcommand{\cf}{cf.\ }
\newcommand{\resp}{respectively\ }
\newcommand{\Subsection}{\subsection}
\providecommand{\nobreakdash}{\nobreak\hbox{-}}
\setlength{\emergencystretch}{2em}
\multlinegap0pt
\usepackage{algorithm}\usepackage{algpseudocode}
\usepackage{amssymb}
\newtheorem{proposition}{Proposition}
\title{\textbf{On the generalization of $g$-circulant MDS matrices}}
\author{Atif Ahmad Khan, Shakir Ali, Bhupendra Singh}
\date{Source: arXiv:2602.10028v1 (CC BY 4.0)}
\begin{document}
\maketitle

\noindent\emph{Source and licence.} \textbf{On the generalization of $g$-circulant MDS matrices}. Version arXiv:2602.10028v1 (CC BY 4.0), distributed by arXiv under the Creative Commons Attribution 4.0 International licence, http://creativecommons.org/licenses/by/4.0/, as recorded in the arXiv licence metadata for this exact version and shown on the abstract page https://arxiv.org/abs/2602.10028v1. Full licence notice: \texttt{licenses/arxiv-2602.10028-CC-BY-4.0.txt}.\par\smallskip

\noindent\textit{Editorial scope.} Counted unit: the article's proposition 6counting1 (source0.tex lines 384--386). The article's complete original proof of it is delivered with it (source0.tex lines 387--413, one \texttt{proof} environment of 27 lines). No mathematics is written, repaired or re-derived: the statement and the proof are the article's own lines, in the article's own notation, and no retained line is substituted or re-flowed.  Retained uncounted as the source-local context the proof needs: lines 317--323.  Nothing was omitted from any retained span, so the package carries no omission record.  No retained line contains a citation, so the wrapper carries no bibliography and no bibliography entry.  No retained line contains a figure environment, an image-inclusion command or drawing code, so no image is delivered and the package declares no asset dependency.  Editorial formatting only, in addition: an AMS \texttt{amsart} preamble that restates the article's own declarations and macros used by the retained lines, namely the environments \texttt{proposition} on separate counters, together with the standard packages those lines need.  The retained environments print in this document as Proposition 1 in order of appearance, whereas the article prints the same items with the numbering of its own full document.  The complete native source of the article as distributed in the arXiv e-print for this exact version is bundled unchanged as \texttt{sources/194367/source0.tex}.
\begin{proposition}\label{6Proposition1}
	Let $h(x)=x^m-\lambda+\sum_{i=0}^{m-1}h_ix^i\in \mathbb{F}_q[x],$ where $\lambda\in\mathbb{F}^*_q$ such that $g\equiv1\mod \textup{ord}(\lambda)$. Then, the map 
	\begin{eqnarray*}
		\phi:\frac{\mathbb{F}_q[x]}{(x^m-\lambda)} &\longrightarrow& \frac{\mathbb{F}_q[x]}{(x^m-\lambda)}\\		\phi(\bar{Q}(x))&:=&\bar{Q}(x^g)\bar{h}(x),\textup{ for all $\bar{Q}(x)$ in \( \mathbb{F}_q[x]/(x^m - \lambda) \),}
	\end{eqnarray*}
	is a vector space linear transformation.
\end{proposition}

\begin{proposition}\label{6counting1}
The upper bound of the total number of consta-$g$-circulant matrix associated with the Proposition \ref{6Proposition1} is $\left( \left\lfloor \frac{m - 1}{\operatorname{ord}(\lambda)} \right\rfloor + 1 \right) \cdot q^m$, where  $\lfloor \cdot \rfloor$ represent the greatest integer function.
\end{proposition}\label{6Proposition002}

\begin{proof}
	Let us denote \( r = \operatorname{ord}(\lambda) \), where \( \lambda \in \mathbb{F}_q^* \), and let \( m \) be a positive integer. According to Proposition~\ref{6Proposition1}, for each integer \( g \in \{1, 2, \dots, m\} \), the map
	\[
	\phi: \frac{\mathbb{F}_q[x]}{(x^m - \lambda)} \to \frac{\mathbb{F}_q[x]}{(x^m - \lambda)}, \quad \phi(\bar{Q}(x)) := \bar{Q}(x^g)\bar{h}(x)
	\]
	is a vector space linear transformation, provided that
	\(
	g \equiv 1 \pmod{r}.
	\) We seek the number of integers \( g \in \{1, 2, \dots, m\} \) such that \( g \equiv 1 \mod r \). These are precisely the values in the arithmetic progression
	\[
	g = 1,\, 1 + r,\, 1 + 2r,\, \dotsc \leq m.
	\]
	Let \( N_g \) be the number of such values. This is the number of integers of the form \( g = 1 + kr \) for integers \( k \geq 0 \), such that
	\[
	1 + kr \leq m \quad \Rightarrow \quad k \leq \left\lfloor \frac{m - 1}{r} \right\rfloor.
	\]
	Thus,
	\[
	N_g = \left\lfloor \frac{m - 1}{r} \right\rfloor + 1.
	\]
	
\noindent 	Note that \( \bar{h}(x) \in \mathbb{F}_q[x]/(x^m - \lambda) \), which is a vector space of dimension \( m \) over \( \mathbb{F}_q \). Therefore, the number of possible choices for \( \bar{h}(x) \) is $q^m.$ Each admissible value of \( g \) yields \( q^m \) possible linear maps \( \phi \), corresponding to the choices of \( \bar{h}(x) \). Hence, the total number of such transformations is atmost
	$$ N_g \cdot q^m = \left( \left\lfloor \frac{m - 1}{r} \right\rfloor + 1 \right) \cdot q^m.
	$$
	Since each linear transformation \( \phi \) corresponds to a consta-$g$-circulant matrix, this is also the upper bound of  such matrices.
	
\end{proof}

\end{document}
